% ============================================================================================
% BAB III ANALISIS MASALAH — KERANGKA (bullet), bukan prosa jadi.
%
% PEMETAAN. DECISIONS.md memakai skema lima bab dan menamai babnya
% "Bab III — Analisis dan Metodologi", berisi tujuh seksi (§1 sampai §7).
% Template TA memakai tujuh bab. Kesepakatan yang dipakai di sini:
%
%   DECISIONS §1 Sources and acquisition   -> Bab III (berkas ini)
%   DECISIONS §2 Quality and exclusions    -> Bab III (berkas ini)
%   DECISIONS §3 Shape of the dataset      -> Bab III (berkas ini)
%   DECISIONS §7 Requirements              -> Bab III, subbab Kebutuhan Nonfungsional
%   DECISIONS §4 Target and preprocessing  -> Bab IV Perancangan
%   DECISIONS §5 Experimental protocol     -> Bab IV Perancangan
%   DECISIONS §6 Implementation            -> Bab IV Perancangan (sebagian ke Pemilihan Solusi)
%   DECISIONS "Bab IV — Hasil"             -> Bab VI Evaluasi
%
% Jadi "Bab IV" pada DECISIONS.md berarti HASIL, sedangkan Bab IV pada dokumen
% ini adalah PERANCANGAN. Jangan membaca medan `Bab` di DECISIONS.md sebagai
% tujuan berkas; ia label topik.
%
% PENYIMPANGAN DARI TEMPLATE (D-45). Template menyediakan tiga subbab
% (Analisis Kondisi Saat Ini, Analisis Kebutuhan, Analisis Pemilihan Solusi) dan
% menyatakan sendiri "Pembagian subbab tidak rigid dan dapat bervariasi". Tiga
% subbab data ditambahkan di antaranya untuk menampung §1--§3. Subbab template
% dipertahankan seluruhnya, beserta \label-nya.
%
% Dua sitasi template (\textcite{laudon2020} dan \autocite{pressman2019}) DIHAPUS.
% daftar-pustaka.bib di paper/Rafly TA/ masih berkas dummy, sehingga kunci itu
% teratasi ke entri dummy tanpa memberi peringatan — kegagalan senyap.
% ============================================================================================
\cleardoublepage
\chapter{ANALISIS MASALAH}
\label{chap:analisis-masalah}

\section{Analisis Kondisi Saat Ini}

\begin{itemize}
    \item Model konseptual saat ini: GFD diestimasi pada resolusi
          \textbf{bulanan}, dari rekaman jaringan deteksi tunggal, untuk satu
          domain iklim.
          % CITE: studi predecessor — pemodelan GFD bulanan Jawa Barat
    \item Tiga keterbatasan yang menjadi titik masuk TA ini:
          \begin{itemize}
              \item resolusi bulanan membuang resolusi yang sebenarnya dimiliki
                    ERA5 dan NASA POWER, yang keduanya per jam;
              \item satu domain iklim tidak memungkinkan generalisasi diukur;
              \item indeks K (\texttt{KX}) yang dipakai studi terdahulu tidak
                    tersedia di kedua domain (D-01), sehingga perbandingan
                    langsung melemah dan hal itu harus dinyatakan.
          \end{itemize}
    \item Perpindahan ke sel-jam (D-21) adalah perubahan kategori, bukan
          sekadar penghalusan: satu baris menjadi satu sel 0,5$^\circ$ untuk
          satu jam, dan target berubah menjadi proses cacah yang 94--97\%
          bernilai nol.
          % TODO(open): Pertanyaan terbuka 7 — argumen positif untuk resolusi
          %   jam HARUS DITULIS PENUH DI SINI, dan belum pernah ditulis di mana
          %   pun. Yang ada di DECISIONS.md hanya recollection: "pemetaan
          %   masukan-keluaran yang lebih akurat antara keadaan meteorologi dan
          %   petir yang menyertainya". Karena resolusi jam menyimpang dari
          %   mitigasi kelangkaan data yang disetujui pada proposal, argumen ini
          %   harus disajikan BERSAMA tiga keberatan yang menyertainya, bukan
          %   menggantikannya. Ini penulisan, bukan pencarian data.

\end{itemize}

\section{Analisis Sumber dan Akuisisi Data}
\label{sec:sumber-data}

% <- DECISIONS.md Bab III §1

\begin{itemize}
    \item \textbf{PLN Puslitbang LDS (tropis, Jawa Barat).} Seluruh
          \textit{sheet} pada setiap \textit{workbook} dibaca, sehingga
          reorganisasi ekspor oleh pemasok tidak merusak pemuat. Baris
          dipertahankan bila \texttt{Discrimination} berawalan \texttt{CG};
          polaritas diturunkan dari sufiks \texttt{+}/\texttt{-} pada medan yang
          sama. Hasil: \textbf{2.242.100 sambaran CG}, 2018-01-01 sampai
          2024-12-31 --- angka pasca-filter, bukan cacah baris mentah (D-22).
    \item \textbf{NASA MERLIN (subtropis, Florida).} KSC Weather Archive tidak
          memiliki API terdokumentasi, sehingga URL ekspor direkonstruksi:
          setiap blok 7 karakter adalah satu \textit{datetime} basis-62.
          \textbf{89 ekspor} berjendela \textbf{29 hari} (bukan 30), dan 89
          jendela cocok dengan 89 berkas secara dua arah. Hasil:
          \textbf{5.301.491 sambaran} (D-23).
    \item \textbf{NASA POWER.} Titik, bukan regional, pada resolusi jam ---
          \textit{endpoint} regional per jam tidak ada dan mengembalikan 404.
          Permintaan mengikuti geometri MERRA-2 (lintang kelipatan
          0,5, bujur kelipatan 0,625), dipotong per tahun
          (\texttt{POWER\_HOURLY\_CHUNK = "Y"}). \textbf{84 titik native}, tropis
          6~lintang $\times$ 5~bujur = 30 (D-24).
    \item \textbf{ERA5.} Satu permintaan CDS per domain per bulan
          (\texttt{ERA5\_HOURLY\_CHUNK = "M"}): $7 \times 12 \times 2 =
          \textbf{168 permintaan}$, \textbf{168 berkas, 258 MB}. Enam peubah
          aras tunggal: CAPE, indeks K, air es dan air cair kolom total, serta
          dua divergensi fluks; seluruh 24 jam diminta. Pemotongan per
          \textbf{tahun} sudah diuji tangan dan \textbf{ditolak} --- satu
          domain-tahun kira-kira 53.000 medan, melewati batas biaya CDS
          (D-25).
\end{itemize}

% CITE: dokumentasi ERA5 / Copernicus CDS
% CITE: dokumentasi NASA POWER
% CITE: KSC Weather Archive, produk MerlinCloudToGround

\section{Analisis Kualitas dan Eksklusi Data}
\label{sec:kualitas-data}

% <- DECISIONS.md Bab III §2

\begin{itemize}
    \item \textbf{Tiga bulan MERLIN tanpa data, dikecualikan dan tidak diisi}
          (D-26). Lima dari 89 ekspor tidak mengembalikan isi yang berguna,
          sehingga \textbf{2021-03, 2022-12, dan 2024-02} tidak hadir sama
          sekali pada tabel terproses --- bukan diisi nol. Akibatnya cakupan
          subtropis \textbf{81 dari 84 bulan}, lawan \textbf{84 dari 84} pada
          tropis. Ini isi yang hilang, bukan berkas yang hilang.
          % TODO(open): Pertanyaan terbuka 12 — apakah 2022-12 dan 2024-02
          %   merupakan pemadaman jaringan atau memang periode tenang? Belum
          %   dipisahkan.
    \item \textbf{Jam pada ekspor PLN adalah waktu lokal Jakarta, divalidasi}
          (D-17). Tidak ada keterangan zona waktu pada ekspor, sehingga ini
          awalnya asumsi dan kini hasil fisis: histogram atas 2.242.100
          sambaran memuncak pada \textbf{16:00 lokal dengan 431.544 flash},
          diapit 15:00 dan 17:00, dengan palung 08:00--10:00. Seandainya ekspor
          bersatuan UTC, puncak lokal sesungguhnya jatuh pada 23:00 --- maksimum
          tengah malam untuk konveksi tropis, yang tidak fisis.
    \item \textbf{Gradien \textit{detection efficiency} subtropis terukur dan
          terkonfound} (D-27). Jumlah \textit{flash} meluruh terhadap jarak dari
          Cape dengan Spearman $-0,955$, melintasi empat orde besaran
          (28,75/$-$80,75: 31 km, 444.193 flash; 26,75/$-$78,75: 267 km, 87
          flash). Dilaporkan sebagai \textbf{konfound antara respons instrumen
          dan klimatologi yang tidak dapat dipisahkan dengan MERLIN saja} ---
          bukan artefak, dan bukan fisika.
\end{itemize}

% TODO(open): Pertanyaan terbuka 10 — KOMPOSISI CG/IC EKSPOR MERLIN.
%   Status: setengah terselesaikan. Tidak ada kolom diskriminasi pada MERLIN.
%   Bukti yang ada: MERLIN 5,7% positif lawan PLN 14,1% pada arus puncak
%   bertanda atas seluruh rekaman; fraksi positif beramplitudo rendah berdekatan
%   (3,1% lawan 2,6%). Kontaminasi IC MENAIKKAN fraksi positif, dan MERLIN
%   berada jauh DI BAWAH rujukan CG-only, jadi arah buktinya mendukung CG-only.
%   Sisi rujukan kini terverifikasi (pertanyaan terbuka 5: ekspor PLN CG-only di
%   sumbernya untuk ketujuh tahun), sehingga hanya satu sisi perbandingan yang
%   tak pasti — inferensinya lebih kuat, tetapi TETAP TIDAK KONKLUSIF.
%   Settled by: dokumentasi KSC Weather Archive untuk produk
%   MerlinCloudToGround. BUKAN dengan analisis lanjutan — data sudah memberi apa
%   yang bisa diberikannya.
%   DIBUTUHKAN SEBELUM BAB III SELESAI, karena menentukan target subtropis
%   disebut apa. Sampai selesai: SETIAP hasil subtropis bersifat PROVISIONAL dan
%   harus dinyatakan demikian, di Bab VI dan di Bab VII.

\section{Analisis Bentuk Dataset}
\label{sec:bentuk-dataset}

% <- DECISIONS.md Bab III §3

\begin{itemize}
    \item \textbf{Resolusi temporal: satu baris adalah satu sel-jam} (D-21).
          \texttt{TIME\_FREQ = "h"}. Jalur kode harian dan bulanan tetap
          berfungsi; tidak ada yang dibangun di keduanya.
    \item \textbf{Grid spasial 0,5$^\circ$, kotak dibulatkan ke luar} (D-15).
          Pusat sel jatuh pada kelipatan 0,5$^\circ$ ditambah 0,25$^\circ$. Luas
          sel dihitung dari kosinus lintang, sehingga sel Florida benar-benar
          lebih kecil daripada sel Jawa Barat. Alasan angka 0,5$^\circ$: cocok
          dengan studi terdahulu \textbf{dan} dengan resolusi lintang native
          NASA POWER --- dua alasan sekaligus. Biayanya: pembulatan ke luar
          itulah yang menghasilkan pelek nol pada tropis.
    \item \textbf{Pelabelan bin: UTC, dipaksa} (D-16). Pada
          \texttt{TIME\_FREQ = "h"} setelan \texttt{TZ\_MODE} diabaikan dan
          dipaksa UTC. Pada resolusi jam persoalannya bukan badai mana masuk bin
          mana --- satu jam adalah jam yang sama pada kedua jam dinding ---
          melainkan bagaimana bin itu \textbf{dilabeli}; keempat sumber harus
          sepakat atau \textit{join} diam-diam menghasilkan tabel kosong. Siklus
          diurnal dibawa terpisah oleh kolom \texttt{hour\_of\_day\_local}.
    \item \textbf{Gerbang cakupan dimatikan: \texttt{MIN\_COVERAGE = 0.0}}
          (D-18). Tidak ada bulan yang dibuang atas dasar cakupan. Gerbang 0,9
          akan menghapus \textbf{56 dari 81 bulan subtropis dan 26 dari 84
          tropis} --- sekitar 70\% Florida --- dan menghapus \textbf{paruh tahun
          yang berlawanan} pada kedua domain (tropis menipis pada musim kemarau
          JJAS, subtropis pada musim dingin DJF). Kedua himpunan latih tidak lagi
          merentang julat musiman yang sebanding, sehingga selisih lintas domain
          yang diukur sesudahnya sebagian menjadi artefak filter.
    \item \textbf{Sel kosong diisi nol eksplisit} (D-19). Membuangnya berarti
          diam-diam melatih model pada "periode yang memang berpetir", yang
          menaikkan rerata target dan membuat R\textsuperscript{2} tampak lebih
          baik daripada sebenarnya.
          \textit{Catatan yang harus dipatuhi:} pembenaran kedua yang pernah
          diingat --- bahwa nol berarti "tidak ada galat sensor" --- \textbf{tidak
          boleh dipakai}, karena D-18 membantahnya dengan pengukuran proyek ini
          sendiri (cakupan hari minimum tropis 12,90\%).
    \item \textbf{Inflasi nol dibawa dan dilaporkan, tidak direkayasa hilang}
          (D-20). Terukur pada \textit{build} 2026-09-06: \textbf{94,30\%
          tropis, 97,00\% subtropis}. Tidak ada penyeimbangan ulang saat
          pembangunan. Dua konsekuensinya adalah kewajiban \textbf{pelaporan},
          bukan masalah pemodelan: R\textsuperscript{2} nyaris tak bermakna
          sehingga pangsa nol dilaporkan di sebelah setiapnya, dan setiap klaim
          harus mengalahkan \textit{baseline} trivial lebih dahulu.
    \item \textbf{Tidak ada nilai hilang} pada prediktor mana pun di kedua
          tabel, sehingga \textbf{tidak ada langkah imputasi} di mana pun dalam
          paket pemodelan. Nyatakan ini secara eksplisit.
\end{itemize}

% Tabel ringkas "build of record" (2026-09-06) yang harus masuk subbab ini —
% angka-angka ini kanonik, ambil dari DECISIONS.md baris 86, jangan tulis ulang
% dari ingatan:
%   Sambaran            2.242.100      / 5.301.491
%   Sel                 30             / 56
%   Sel-jam             1.841.040      / 3.314.304
%   Sel-jam tak nol     104.955        / 99.273
%   Pangsa nol          94,30%         / 97,00%
%   Julat tak nol       1..2.391       / 1..11.777
%   Bulan tercakup      84 dari 84     / 81 dari 84
%   Cakupan hari, rerata (min)  85% (13%) / 59% (3%)   <- BERBOBOT BARIS
%   Prediktor           14, KX ada     / 13, KX absen
%   Nilai hilang        nihil          / nihil
% CATATAN: nyatakan pembobotannya. Versi lama tabel ini menulis tropis 86%,
% yang mendahului pengukuran dan kemungkinan berbobot bulan. Angka berbobot
% baris: tropis rerata 0,8549 min 0,1290; subtropis rerata 0,5921 min 0,0323.
% Angka cakupan berasal dari kolom `coverage` pada parquet terbangun; .meta.json
% TIDAK memuatnya.

\section{Analisis Kebutuhan}

\subsection{Identifikasi Masalah Pengguna}

\begin{itemize}
    \item Pengguna: perencana proteksi sistem tenaga listrik yang memerlukan
          estimasi kerapatan sambaran pada wilayah dan periode tertentu.
    \item Masalah 1: estimasi bulanan tidak menjawab pertanyaan pada skala
          kejadian konvektif.
    \item Masalah 2: estimasi yang dilatih pada satu domain iklim belum
          diketahui keteralihannya ke domain lain.
    \item Masalah 3: jaringan deteksi memiliki cakupan tidak seragam, dan
          angka yang tidak menyertakan cakupannya tidak dapat dinilai
          pembacanya (D-10).
\end{itemize}

\subsection{Kebutuhan Fungsional}

\begin{itemize}
    \item \textbf{F-01.} Sistem membangun tabel prediktor--target beresolusi
          sel-jam dari empat sumber untuk kedua domain, 2018--2024.
    \item \textbf{F-02.} Sistem melatih dua \textit{arm} model (kuantum dan
          klasik) dalam satu \textit{loop} latih yang sama.
    \item \textbf{F-03.} Sistem melaporkan metrik pada empat jendela pelaporan
          1/3/6/24 jam (D-04).
    \item \textbf{F-04.} Sistem mengevaluasi model lintas domain dengan
          \textit{scaler} domain sumber (D-36).
    \item \textbf{F-05.} Setiap hasil membawa \textit{snapshot} konfigurasi yang
          memproduksinya (D-42).
          % TODO(open): hasil sebelum 2026-09-08 TIDAK memuat `feature_reduction`
          %   dan `n_qubits` pada snapshot-nya. Periksa provenance sebelum
          %   mengutip metrik mana pun.
\end{itemize}

\subsection{Kebutuhan Nonfungsional}

% <- DECISIONS.md Bab III §7

\begin{itemize}
    \item \textbf{NF-01: NRMSE $< 0,1$ dan R\textsuperscript{2} $> 0,9$.}
          Dipertahankan \textbf{apa adanya}, dan dinyatakan di muka
          \textbf{tidak akan terpenuhi} pada resolusi jam (D-12). Alasannya
          perbedaan kategori, bukan kegagalan pemodelan: NF-01 ditulis terhadap
          kerapatan GFD bulanan, sementara target pada resolusi jam adalah
          proses cacah 94--97\% nol, dan R\textsuperscript{2} di atasnya nyaris
          tak bermakna. NRMSE dinormalkan terhadap \textbf{julat} teramati ---
          konvensi yang dipakai saat NF-01 ditulis --- sehingga pada target
          sejarang ini ia didominasi nilai maksimum, dan dilaporkan
          \textbf{di samping} RMSE, bukan menggantikannya.
    \item \textbf{Apa yang diubah D-04, dan seberapa besar: lebih kecil dari
          harapan.} Pada 24 jam pangsa nol tropis turun ke 64,63\%, sehingga
          R\textsuperscript{2} menjadi dapat ditafsirkan, tetapi target tetap
          didominasi nol. NF-01 dilaporkan terhadap \textbf{setiap} jendela pada
          kurva D-04 --- jawaban yang jauh lebih kuat daripada meleset diam-diam
          atau mendefinisikan ulang diam-diam.
          % TODO(open): status NF-01. D-12 memerintahkan hal ini diangkat dengan
          %   pembimbing SEBELUM Bab VII (Penutup), bukan di dalamnya. BELUM
          %   DILAKUKAN. Yang sudah pasti: keputusannya adalah melaporkan
          %   tidak-terpenuhi dan menjelaskannya. Yang belum: persetujuan
          %   pembimbing atas cara pelaporan itu.
    \item \textbf{NF-02: reprodusibilitas}, diperlakukan sebagai kebutuhan
          \textbf{bukti} dengan artefak bernama, bukan sebagai satu baris pada
          tabel (D-42). Empat mekanisme membawanya:
          \begin{itemize}
              \item satu \texttt{RANDOM\_SEED = 18222067} didefinisikan sekali
                    pada konfigurasi pipeline dan \textbf{diimpor, bukan
                    disalin} oleh konfigurasi pemodelan, sehingga "seed
                    18222067" tidak dapat berarti dua hal; seed per-run
                    $(0,1,2)$ terpisah dan merupakan sumbu replikasi (D-29);
              \item \textit{snapshot} konfigurasi menyertai setiap hasil
                    (\texttt{.meta.json} dan \texttt{\_config\_snapshot()});
              \item satu lingkungan tersemat eksak di akar repositori
                    (D-41), diverifikasi pada Python 3.13.7 dengan 3.10 sebagai
                    lantai;
              \item \texttt{data/processed/} dapat dibangun ulang dengan satu
                    perintah.
          \end{itemize}
          % TODO(open): NF-02 terpenuhi SECARA MEKANISME dan BELUM DIUJI
          %   ujung-ke-ujung (D-42; pertanyaan terbuka 28 — replay penuh belum
          %   pernah dijalankan). Tulis persis kalimat itu, jangan lebih.
    \item \textbf{Kendala redistribusi data PLN} (D-43): rekaman mentah datang
          melalui pembimbing, tidak publik, dan tidak boleh diredistribusi.
          Ini membatasi bentuk bukti reprodusibilitas yang dapat ditawarkan
          kepada pihak luar, dan harus dinyatakan bersama NF-02.
\end{itemize}

\section{Analisis Pemilihan Solusi}

\subsection{Alternatif Solusi}

\begin{itemize}
    \item \textbf{Alternatif 1 --- regresi cacah model tunggal.} Ditolak sebagai
          bawaan: kurva jendela terukur pada D-04 mematahkan asumsi bahwa
          jendela pelaporan yang lebar akan membuat regresi cacah model tunggal
          terkondisikan baik. Pada 6 jam target masih 85,40\% / 91,35\% nol
          (D-13).
    \item \textbf{Alternatif 2 --- \textit{hurdle model} dua tahap.}
          \textbf{Dipilih sebagai bawaan} (\texttt{TASK = "hurdle"}): tahap
          okurensi biner, lalu regresi \texttt{log1p(flash\_count)} pada baris
          tak nol saja (D-13). Rinciannya di Bab~\ref{chap:perancangan}.
    \item \textbf{Alternatif 3 --- pustaka pemodelan siap pakai
          (\texttt{MLPRegressor}).} Ditolak: estimator scikit-learn memiliki
          \textit{loop} latihnya sendiri, sehingga optimizer, ukuran
          \textit{batch}, aturan \textit{early stopping}, dan penanganan seed
          tidak dapat disamakan. Setiap perbandingan yang ditarik terhadapnya
          membawa konfound, dan "NN-nya dilatih dengan cara berbeda" adalah
          keberatan pertama yang akan diajukan penguji (D-09).
    \item \textbf{Alternatif 4 --- Qiskit sebagai jalur gradien produksi.}
          Ditolak atas dasar biaya terukur (D-28). Lihat subbab berikutnya.
\end{itemize}

\subsection{Analisis Penentuan Solusi}

\begin{itemize}
    \item \textbf{Kedua model dilatih dalam satu \textit{loop} PyTorch}
          (D-09): optimizer, ukuran \textit{batch}, objek \textit{loss}, aturan
          \textit{early stopping}, dan penanganan seed yang sama. Satu-satunya
          perbedaan antara kedua model harus \textbf{lapisannya}.
    \item \textbf{Gradien dihitung PennyLane \texttt{lightning.qubit} dengan
          metode \texttt{adjoint}, bukan Qiskit} (D-28). Sirkuit tetap
          didefinisikan dan diverifikasi terhadap Qiskit; jalur
          \textit{parameter-shift} milik \texttt{qiskit-machine-learning} tidak
          dipakai untuk \textit{run} produksi. Dasarnya adalah tolok ukur per
          sampel yang terukur, bukan preferensi.
          % CITE: PennyLane
          % CITE: Qiskit
          % CITE: adjoint differentiation
          % TODO(open): Pertanyaan terbuka 9 — konsekuensi keputusan ini bagi
          %   frasa "Framework Qiskit" pada judul. Satu percakapan pembimbing,
          %   digabung dengan pertanyaan 8.
    \item \textbf{Bukti bahwa kedua jalur gradien menghasilkan model yang sama}
          direncanakan sebagai satu \textit{run} Qiskit tereduksi (D-39):
          \textbf{500 baris $\times$ 15 epoch, satu fold, satu seed, tahap
          okurensi saja} --- sekitar 55 menit. Pasangan \textit{hurdle} penuh
          pada ukuran penuh adalah 21,9 jam dan \textbf{tidak dijalankan}.
          Alasannya: \texttt{verify\_against\_qiskit()} membuktikan kesepakatan
          \textbf{forward pass} sampai 1e-10, dan itu menjawab pertanyaan
          tentang \textit{sirkuit}; pertanyaan di sidang adalah tentang
          \textit{pelatihan}.
          % TODO(open): D-39 berstatus DECIDED, NOT YET RUN (pertanyaan terbuka
          %   24). Jangan menuliskannya seolah sudah dijalankan.
    \item \textbf{Batas yang dinyatakan di muka, bukan ditemukan penguji:}
          \begin{itemize}
              \item seluruh hasil subtropis provisional sampai pertanyaan
                    terbuka 10 selesai;
              \item jendela 6 jam adalah \textbf{nominasi demi keterbacaan},
                    bukan optimum yang diturunkan --- tidak ada kriteria seleksi
                    yang pernah diterapkan padanya (D-04);
              \item \textit{sweep} memakai satu fold dan satu seed, sementara
                    tabel utama memakai tiga (D-29, D-37);
              \item tujuh hiperparameter tidak memiliki alasan tercatat
                    (pertanyaan terbuka 22). Ini keterbatasan yang dinyatakan,
                    \textbf{bukan lubang yang diisi} dengan pembenaran yang
                    ditulis belakangan.
          \end{itemize}
\end{itemize}
